% SEGMENT OLP-0576-S01
% Part: set-theory
% Chapter: ord-arithmetic
% Section: addition

\documentclass[../../../include/open-logic-section]{subfiles}

\begin{document}

\olfileid{sth}{ord-arithmetic}{add}
\olsection{வரிசையெண் கூட்டல்}

$\alpha$ வையும் $\beta$ வையும் கூட்ட விரும்புகிறோம் என்று கொள்வோம்.
$\alpha$ வின் ஒரு பிரதிக்குப் பின்னால் $\beta$ வின் ஒரு பிரதியை
உடனடியாக வைக்கலாம். (\olref[ordinals][basic]{ordinalsaresubsets} இன்
படி $\alpha \subseteq \beta$ அல்லது $\beta \subseteq \alpha$
என்பதால், அவற்றின் \emph{பிரதிகளை} எடுக்க வேண்டும்.) உள்ளுணர்வாக,
முதலில் $\alpha$-தொடர் படிகளைக் கடந்து, \emph{அதற்குப் பிறகு}
$\beta$-தொடர் படிகளைக் கடப்பதாக இது அமையும். கிடைக்கும் தொடர் நன்கு
வரிசைப்படுத்தப்பட்டிருக்கும்; ஆகவே
\olref[ordinals][ordtype]{thmOrdinalRepresentation} இன் படி,
அது (ஒரேயொரு) வரிசையெண்ணுடன் வரிசை ஓரினச்சார்பு உடையதாகும். அந்த
வரிசையெண்ணை $\alpha$, $\beta$ ஆகியவற்றின் \emph{கூட்டுத்தொகை} எனக்
கருதலாம்.

வரிசையெண் கூட்டலின் உள்ளுணர்வுக் கருத்து இதுவே. இதைத் துல்லியமாக
வரையறுக்க, கணங்களின் \emph{பிரதிகளை} உருவாக்கும் கருத்திலிருந்து
தொடங்குகிறோம். எந்தப் பொருள் எங்கிருந்து வந்தது என்பதைப் பதிவு செய்ய,
தன்னிச்சையான $0$, $1$ குறிச்சொற்களைப் பயன்படுத்துகிறோம்:

\begin{defn}\ollabel{defdissum}
$A$, $B$ ஆகியவற்றின் \emph{வெட்டாக் கூட்டுத்தொகை} என்பது
$A \disjointsum B = (A\times \{0\}) \cup (B \times \{1\})$.
\end{defn}

% SEGMENT OLP-0576-S02
அடுத்து, வரிசையெண் இணைகளின் மீது ஒரு வரிசையை வரையறுக்கிறோம்:

\begin{defn}
எந்த வரிசையெண்கள் $\alpha_1, \alpha_2, \beta_1, \beta_2$ க்கும்:
\begin{align*}
  \tuple{\alpha_1, \alpha_2} \rlexless \tuple{\beta_1, \beta_2}\text{ எனில், எனில் மட்டுமே }&
  \text{$\alpha_2 \in \beta_2$ ஆகும்; அல்லது}\\
  &\text{$\alpha_2 = \beta_2$ மற்றும் $\alpha_1 \in \beta_1$ ஆகிய இரண்டும் ஆகும்.}
\end{align*}
\end{defn}

இது \emph{தலைகீழ் அகராதி வரிசை}: முதலில் இரண்டாவது
உறுப்பின்படியும், அதன் பிறகு முதலாவது உறுப்பின்படியும் வரிசைப்படுத்துகிறோம்.
$\alpha$ வின் பிரதிக்குப் பின்னால் $\beta$ வின் பிரதியை வைத்தால்
கிடைக்கும் வரிசை வகையாக $\alpha \ordplus \beta$ வைப் வரையறுக்க
விரும்பினோம் என்பதை நினைவுகூருங்கள். அதற்காகப் பின்வருமாறு கூறுகிறோம்:

\begin{defn}\ollabel{defordplus}
எந்த வரிசையெண்கள் $\alpha$, $\beta$ க்கும் அவற்றின் கூட்டுத்தொகை
$\alpha \ordplus \beta = \ordtype{\alpha \disjointsum \beta, \rlexless}$.
\end{defn}
\noindent
இங்கு குறியீட்டைச் சற்றுத் தளர்வாகப் பயன்படுத்தியுள்ளோம். கண்டிப்பாகச்
சொன்னால், ``$\rlexless$'' என்பதற்குப் பதிலாக
``$\Setabs{\tuple{x,y}\in \alpha\disjointsum\beta}{x \rlexless y}$''
என்று எழுத வேண்டும். சுருக்கத்திற்காகத் தொடர்ந்து இதே குறியீட்டைத்
தளர்வாகப் பயன்படுத்துவோம்.

பின்வரும் முடிவும் \olref[ordinals][ordtype]{thmOrdinalRepresentation}
உம் சேர்ந்து, நமது வரையறை முறையாக அமைந்திருப்பதை உறுதிசெய்கின்றன:

\begin{lem}\ollabel{ordsumlessiswo}
எந்த வரிசையெண்கள் $\alpha$, $\beta$ க்கும்
$\tuple{\alpha \disjointsum \beta, \rlexless}$ ஒரு நன்கு
வரிசைப்படுத்தலாகும்.
\end{lem}

\begin{proof}
$\alpha \disjointsum \beta$ இன் மீது $\rlexless$ ஒப்பிடத்தக்கது
என்பது தெளிவு. ஒவ்வொரு வெறுமையற்ற உட்கணத்திற்கும் மீச்சிறு உறுப்பு
உள்ளது என்பதைக் காட்ட, வெறுமையல்லாத
$X \subseteq \alpha \disjointsum \beta$ ஒன்றை எடுத்துக்கொள்வோம்.
இரண்டாவது ஆயம் இயன்ற அளவு சிறியதாக உள்ள $X$ இன் உறுப்புகளின்
உட்கணத்தை $Y$ என்க; அதாவது,
$Y = \Setabs{\tuple{\gamma, i} \in X}{(\forall \tuple{\delta, j} \in X)i \leq j}$.
இப்போது $Y$ இல் முதலாவது ஆயம் மிகச் சிறியதாக உள்ள உறுப்பைத்
தேர்ந்தெடுக்கலாம்.
\end{proof}
\noindent
இதனால் வரிசையெண் கூட்டலுக்கு ஒரு தெளிவான, வெளிப்படையான வரையறை
கிடைத்தது. ஆச்சரியமில்லாத ஒரு உண்மை இதோ
(\olref[z][infinity-again]{defnomega} இன் படி $1 = \{0\}$ என்பதை
நினைவுகூருங்கள்):
\begin{prop}
எந்த வரிசையெண் $\alpha$ க்கும்
$\alpha \ordplus  1 = \ordsucc{\alpha}$.
\end{prop}

\begin{proof}
% SOURCE CORRECTION TA-STH-023: the two already tagged components form a union.
$\ordsucc{\alpha} = \alpha \cup \{\alpha\}$ இலிருந்து
$\alpha\disjointsum1 = (\alpha \times \{0\})
\cup (\{0\} \times \{1\})$ க்கான வரிசை ஓரினச்சார்பு $f$ ஐக்
கவனியுங்கள்: $\gamma \in \alpha$ எனில்
$f(\gamma) = \tuple{\gamma, 0}$; மேலும்
$f(\alpha) = \tuple{0, 1}$.
\end{proof}
\noindent

% SEGMENT OLP-0576-S03
கூட்டல் சில மீள்வரையறை நிபந்தனைகளைப் பூர்த்திசெய்கிறது என்பதையும்
எளிதில் காட்டலாம்:

\begin{lem}\ollabel{ordadditionrecursion}
எந்த வரிசையெண்கள் $\alpha, \beta$ க்கும்:
\begin{align*}
  \alpha\ordplus 0 &= \alpha\\
  \alpha \ordplus  (\beta\ordplus 1) &= (\alpha \ordplus  \beta) \ordplus  1\\
  \alpha  \ordplus  \beta &= \supstrict_{\delta < \beta}(\alpha \ordplus  \delta) && \text{$\beta$ ஓர் எல்லை வரிசையெண் எனில்}
\end{align*}
\end{lem}

\begin{proof}
ஒவ்வொரு நேர்வையும் தனித்தனியாகச் சரிபார்ப்போம். முதலில்:
% SOURCE CORRECTION TA-STH-024: 0 times {1} is empty, not {0}.
\begin{align*}
  \alpha \ordplus  0
  & = \ordtype{(\alpha \times \{0\}) \cup (0 \times \{1\}), \rlexless} \\
  &= \ordtype{\alpha \times \{0\}, \rlexless}\\
  &= \alpha\\
  \alpha \ordplus (\beta \ordplus  1)
  %&= \ordtype{(\alpha\times \{0\}) \cup (\ordtype{\beta \ordplus  1}\times \{1\}), \rlexless} \\
  &= \ordtype{(\alpha\times \{0\}) \cup (\ordsucc{\beta}\times \{1\}), \rlexless} \\
  &= \ordtype{(\alpha\times \{0\}) \cup (\beta \times \{1\}), \rlexless} \ordplus 1\\
  &= (\alpha \ordplus  \beta) \ordplus  1
\end{align*}
இப்போது $\beta \neq \emptyset$ ஓர் எல்லை வரிசையெண் என்க.
$\delta < \beta$ எனில் $\delta\ordplus 1 < \beta$ உம் ஆகும்;
ஆகவே $\alpha \ordplus  \delta$ என்பது $\alpha \ordplus  \beta$
வின் தகு தொடக்கத் துண்டாகும். எனவே
$X = \Setabs{\alpha \ordplus  \delta}{\delta < \beta}$ இன்
கண்டிப்பான மேல்வரம்பு $\alpha \ordplus  \beta$ ஆகும். மேலும்
$\alpha \leq \gamma < \alpha \ordplus  \beta$ எனில், ஏதோ ஒரு
$\delta < \beta$ க்கு $\gamma = \alpha \ordplus \delta$
என்பது தெளிவு. ஆகவே
$\alpha \ordplus  \beta =
\supstrict_{\delta< \beta}(\alpha\ordplus \delta)$.
\end{proof}

ஆனால் கவனிக்கத்தக்க ஓர் உண்மை உள்ளது. வரிசையெண் கூட்டலை வரையறுக்க,
இதற்குப் \emph{பதிலாக} முடிவிலிக்கடந்த மீள்வரையறைத் தேற்றத்தைப்
பயன்படுத்தி, \olref{ordadditionrecursion} இல் கொடுக்கப்பட்டுள்ள
மீள்வரையறைச் சமன்பாடுகளை மட்டும் விதித்திருக்கலாம்
(அங்கே ``$\beta \ordplus 1$'' என்பதற்குப் பதிலாக
``$\ordsucc{\beta}$'' என்பதைப் பயன்படுத்தி).

ஆக, வரிசையெண்களின் மீது செயல்களை வரையறுக்க இரண்டு வெவ்வேறு வழிகள்
உள்ளன. ஒரு நன்கு வரிசைப்படுத்தப்பட்ட கணத்தை வெளிப்படையாகக்
கட்டமைத்து, அதன் வரிசை வகையை எடுத்துக்கொள்வதன் மூலம் அவற்றை
\emph{கட்டமைப்பு வழியில்} வரையறுக்கலாம். அல்லது மீள்வரையறைச்
சமன்பாடுகளை மட்டும் விதித்து \emph{மீள்வரையறை வழியில்}
வரையறுக்கலாம். ஆனால் சரியாகச் செய்தால், இரண்டிலும் கிடைப்பது ஒன்றே.
ஏனெனில், அந்த மீள்வரையறைச் சமன்பாடுகள் \emph{ஒரேயொரு}
வரிசையெண்ணை நிர்ணயிக்கின்றன என்பதை
\olref[ordinals][ordtype]{thmOrdinalRepresentation} உறுதிசெய்கிறது.

% SEGMENT OLP-0576-S04
பல விதங்களில் வரிசையெண் எண்கணிதம் இயல் எண்களின் கூட்டலைப் போலவே
செயல்படுகிறது. எடுத்துக்காட்டாக, பின்வருவதை நிறுவலாம்:

\begin{lem}\ollabel{ordinaladditionisnice}
$\alpha, \beta, \gamma$ வரிசையெண்கள் எனில்:
\begin{enumerate}
  \item\ollabel{ordaddition1} $\beta < \gamma$ எனில்
  $\alpha \ordplus  \beta < \alpha \ordplus  \gamma$
  \item\ollabel{ordaddition2}
  $\alpha \ordplus  \beta = \alpha\ordplus \gamma$ எனில்
  $\beta = \gamma$
  \item\ollabel{ordaddition3}
  $\alpha \ordplus  (\beta \ordplus \gamma) =
  (\alpha \ordplus  \beta) \ordplus  \gamma$; அதாவது,
  கூட்டல் சேர்ப்புப் பண்புடையது
  \item\ollabel{ordaddition4} $\alpha \leq \beta$ எனில்
  $\alpha \ordplus  \gamma \leq \beta \ordplus \gamma$
\end{enumerate}
\end{lem}

\begin{proof}
மற்றவற்றைப் பயிற்சியாக விட்டுவிட்டு \olref{ordaddition3} ஐ
நிறுவுவோம். \olref{ordadditionrecursion} ஐப் பயன்படுத்தி
$\gamma$ இன் மீது எளிய முடிவிலிக்கடந்த தொகுத்தறிதலால் நிறுவுகிறோம்.
$\gamma = 0$ எனில்:
\[
(\alpha \ordplus  \beta) \ordplus  0 = \alpha \ordplus  \beta  = \alpha \ordplus  (\beta \ordplus  0)
\]
$\gamma = \delta\ordplus 1$ எனில்,
$(\alpha \ordplus  \beta) \ordplus  \delta =
\alpha \ordplus  (\beta \ordplus \delta)$ என்று தொகுத்தறிதல்
எடுகோளாகக் கொள்வோம். இப்போது \olref{ordadditionrecursion} ஐ
மூன்று முறை பயன்படுத்தினால்:
\begin{align*}
  (\alpha \ordplus  \beta) \ordplus  (\delta \ordplus  1) & = ((\alpha \ordplus  \beta) \ordplus  \delta)\ordplus 1\\
  & = (\alpha \ordplus  (\beta \ordplus  \delta)) \ordplus  1\\
  & = \alpha \ordplus  ((\beta \ordplus  \delta)\ordplus 1)\\
  & = \alpha \ordplus  (\beta \ordplus  (\delta\ordplus 1))
\end{align*}
$\gamma$ ஓர் எல்லை வரிசையெண் எனில், ஒவ்வொரு
$\delta \in \gamma$ க்கும்
$(\alpha \ordplus  \beta) \ordplus  \delta =
\alpha \ordplus  (\beta \ordplus  \delta)$ என்று
தொகுத்தறிதல் எடுகோளாகக் கொள்வோம். இப்போது:
\begin{align*}
  (\alpha \ordplus  \beta) \ordplus  \gamma & = \supstrict_{\delta < \gamma}((\alpha \ordplus  \beta) \ordplus  \delta) \\
  &= \supstrict_{\delta < \gamma}(\alpha \ordplus  (\beta \ordplus  \delta))\\
  &= \alpha \ordplus  \supstrict_{\delta < \gamma}(\beta \ordplus  \delta)\\
  &= \alpha \ordplus  (\beta \ordplus  \gamma)
\end{align*}
\end{proof}

\begin{prob}
\olref[sth][ord-arithmetic][add]{ordinaladditionisnice} இன்
மீதமுள்ள பகுதிகளை நிறுவுக.
\end{prob}

% SEGMENT OLP-0576-S05
இவ்விதங்களில் வரிசையெண் கூட்டல் நமக்குப் பழக்கமானதாக இருக்க வேண்டும்.
ஆனால் ஓர் இன்றியமையாத விதத்தில், அது இயல் எண்களின் கூட்டலைப்
\emph{போலல்ல}.

\begin{prop}\ollabel{ordsumnotcommute}
வரிசையெண் கூட்டல் {பரிமாற்றுப் பண்புடையதல்ல};
$1 \ordplus  \omega = \omega < \omega \ordplus  1$.
\end{prop}

\begin{proof}
$1 \ordplus  \omega = \supstrict_{n < \omega} (1 \ordplus n)
= \omega \in \omega \cup \{\omega\} = \ordsucc{\omega} =
\omega \ordplus  1$ என்பதைக் கவனியுங்கள்.
\end{proof}
\noindent
இது முதலில் வியப்பாகத் தோன்றலாம்; ஆனால் \emph{அப்படி
இருக்க வேண்டியதில்லை}. $1 \ordplus  \omega$ வைக் கருதும்போது,
எல்லா இயல் எண்களுக்கும் \emph{முன்னால்} ஓர் உறுப்பைக் கூடுதலாக
வைத்துக் கிடைக்கும் வரிசை வகையைப் பற்றிச் சிந்திக்கிறோம்.
\olref[sfr][infinite][hilbert]{sec} இல் ஹில்பர்ட்டின் விடுதியைப்
பற்றிச் சிந்தித்தது போலச் சிந்தித்தால், இந்தக் கூடுதல் முதலுறுப்பு
ஒட்டுமொத்த வரிசை வகையில் எந்த வேறுபாட்டையும் ஏற்படுத்தாது என்பதை
உள்ளுணரலாம். மறுபுறம், $\omega \ordplus  1$ வைக் கருதும்போது,
எல்லா இயல் எண்களுக்கும் \emph{பின்னால்} ஓர் உறுப்பைக் கூடுதலாக
வைத்துக் கிடைக்கும் வரிசை வகையைப் பற்றிச் சிந்திக்கிறோம்.
அது முற்றிலும் வேறுபட்ட வரிசை வகை!

\end{document}
